\begin{papercard}{knabenhans2026}
\cf{Stages} \stagetag{H} (a worked harm-tree methodology: harm $\leftarrow$ capabilities $\leftarrow$ system components, with the origin of each component), \stagetag{S} (definition of harm and baseline), \stagetag{R} (how to present trees). A 2026 security/privacy paper (13 pages plus appendices, EPFL/ETH/CISPA/MPI-SP) on the European Digital Identity Framework (EUDIF). It uses a formal cryptographic model to show which harms remain even for an ideal privacy-preserving implementation. The subject is digital identity, not health apps; the \emph{harm-tree method} and its use of Citron and Solove's taxonomy are the reusable parts.

\cf{Claim} PETs are often presented as the answer to concerns about digital identity, but the authors ``take a step back and ask whether privacy and PETs are indeed an appropriate answer'' \pc{knabenhans2026}{1}{In this paper, we take a step back and ask whether privacy and PETs are indeed an appropriate answer to concerns about the EUDIF’s negative societal impacts}. Result: ``even a perfectly privacy-preserving implementation of the EUDIF --- one that leaks the minimal information to achieve the use case functionality --- cannot prevent the harms that we identify'' \pc{knabenhans2026}{2}{even a perfectly privacy-preserving implementation of the EUDIF — one that leaks the minimal information to achieve the use case functionality — cannot prevent the harms that we identify}; harms come from the core functionality, from the digital context and from implementation, and only the last class can be addressed by PETs. Three fundamental harm-enabling capabilities (distinguishing, tracking, inference) are shown with cryptographic games.

\cf{Contributions} (1) A formal ``ideal functionality'' $\mathcal{F}_{\text{EUDIF}}$ capturing issuers, users (wallet), relying parties (RPs), issuance, selective disclosure, disclosure policies, RP registration, derived from the eIDAS 2.0 regulation, an implementing regulation and the Architecture and Reference Framework, plus models of age verification and mobile driving licence \pc{knabenhans2026}{1}{We leverage techniques from cryptographic modeling to formally capture the necessary leakage of the functionality of the EUDIF and its proposed applications}. (2) A \emph{harm analysis methodology using harm trees} \pc{knabenhans2026}{1}{Then, we develop a harm analysis methodology that illustrates, using harm trees, how this leakage — and other constraints stemming from design decisions or the context of deployment — lead to harms}. (3) Games for the three capabilities and an argument that regulation, not technology, must restrict the functionality.

\cf{Harm-tree method (Section 4.1)} Built on De and Le M\'etayer's privacy harm analysis. \emph{Three element kinds}: system components, capabilities, harms \pc{knabenhans2026}{6}{We build trees from three kinds of elements: system components, capabilities, and harms}. \emph{System components}: information leakage, outcomes of interactions, requirements imposed on entities (e.g., needing an EUDIF-enabled device). \emph{Capabilities}: what components jointly enable or prevent (track users, inability of a provider to serve a region, inability of a user to access a service). \emph{Harms}: definition ``an impairment or disruption of an entity's capacity or ability to function and exist as it otherwise would have in its usual context'' (CyberPeace Institute), with the current EU practice as the baseline ``usual context''; impairments classified by Citron and Solove and the CyberPeace methodology \pc{knabenhans2026}{6}{taxonomy of privacy harms of Citron and Solove [11] and the CyberPeace Insti- tute’s harms methodology [14]}. Structure: top = harm, second level = capabilities, leaves = components, combined with AND / OR and ``OR+'' (more than one component may be needed) \pc{knabenhans2026}{6}{A harm tree consists of three levels. At the top of the tree is the harm. On the second level, we have the capabilities that enable this harm. At the leaves, we have system components}. \emph{Harm origins} colour each leaf: black = inherent to the core functionality (cannot be removed without stopping the system), blue = inherent to use in a digital context, red = arises from a concrete implementation (the only class PETs can remove). Non-completeness is acknowledged: more harms or pathways would only strengthen the conclusion \pc{knabenhans2026}{7}{We do not claim completeness of the harms, capabilities, or system components we identify in our analysis}.

\cf{Applications} \emph{Age verification}: harms taken from the literature (EDRi, EFF, ...) and Citron and Solove; seven found: discrimination, exacerbated structural exclusion, targeted-advertising harms, reduced availability of resources, censorship, diminished Internet safety, increased cybersecurity risk \pc{knabenhans2026}{7}{We find seven harms: discrimination, exacerbation of structural exclusion, targeted advertising harms, reduced availability of resources, censorship, diminished Internet safety, and increased cybersecurity risk}. Common sub-trees: ``service provider stops service'' (regulatory requirement + cost of compliance) and ``tracking/profiling'' (age information revealed, repeated observation, other attributes, persistent identifier, online metadata). \emph{Mobile driving licence for car rental}: further harms (Section 4.3, Figures 9--12). \emph{Deployment requirements (4.4)}: device binding, certificate checks with registries, disclosure-policy lookups and revocation managers add new parties and flows (``registrar learns $R$'', ``policy manager learns $R$'') that create new tracking pathways, and device requirements increase discrimination and exclusion risk.

\cf{Role of PETs (4.5)} PETs ``can only eliminate a limited subset of pathway to harms''; in every tree ``there exists at least one pathway to the harm that does not require red components'' \pc{knabenhans2026}{10}{in each harm tree, there exists at least one pathway to the harm that does not require red components}. Not all PETs reduce harm equally (batch issuance, pseudonyms); PETs can introduce or exacerbate pathways (opacity leading to non-use or over-disclosure, higher device requirements, concentration of expertise) \pc{knabenhans2026}{10}{PETs can introduce or exacerbate pathways to harm}; for identity-based use cases selective disclosure is not applicable.

\cf{Capabilities by games (Section 5)} \emph{Distinguishing}: ``the ability of a relying party to distinguish between classes of users'' based on one presentation, via differing issuer policies, attribute values, or user behaviour (consent and registration checks) \pc{knabenhans2026}{11}{Informally, distinguishability is the ability of a relying party to distinguish between classes of users}; \emph{trackability}: recognising a user seen before, multi-show or via issuer collusion; \emph{inference}: recovering attributes, including correlated ones, e.g.\ through adaptive requests via several corrupt RPs, unavoidable under an ``any-RP-may-ask-for-anything'' policy \pc{knabenhans2026}{12}{Inference is the adversary’s capability to infer user attributes}. Only regulatory constraints (limits on policies, requestable functions/attributes, no sharing between RPs) can reduce them.

\cf{Takeaway} PETs reduce non-essential leakage, but the intended leakage suffices for many harms; ``promoting PETs as a suitable solution to mitigate harms might result in a false sense of safety'' \pc{knabenhans2026}{13}{promoting PETs as a suitable solution to mitigate harms might result in a false sense of safety}. Appendix A adapts Citron and Solove's categories (economic, physical, psychological, autonomy, discrimination, reputational, relationship) to digital identity with examples \pc{knabenhans2026}{14}{Economic Harms involve financial losses or loss in the value of something}.

\cf{Limits (card)} Domain is digital identity; harms and components are derived from documents and literature by the authors, not from affected people; no ranking, probability or severity; completeness explicitly not claimed; no reliability check on how different analysts would build the trees; the formal model assumes secure channels and ideal functionalities; EU-specific; a 2026 publication without a replication.

\cf{Use for this thesis} (1) A ready, citable structure for the thesis' harm tree (stage H): top node = harm (typed by Citron and Solove), second level = capabilities (e.g., tracking, profiling, inference of health status, re-identification, distinguishing users, disclosure to third parties), leaves = observable \emph{system components} from the walkthrough, policy and store label (SDK present, advertising ID, default-on consent, vague purpose statement). (2) The three-colour origin scheme adapted to apps: inherent to the app's function (e.g., storing cycle data), inherent to the mobile context (OS identifiers, store distribution), and implementation/policy choices (third-party SDKs, consent design); this separates what a better policy could fix from what the product concept implies. (3) The explicit baseline (``usual context'') and ``not claiming completeness'' statement are good practice to copy. (4) A caution against equating label-level privacy measures with harm reduction. (5) The thesis can only observe a subset of components (no network or code analysis), so leaves should be tagged \emph{observed} versus \emph{inferred}.
\end{papercard}
